% SHARD-ID: RC-04Q-H-THETA-REFUTED
% SHARD-TITLE: The Phantasm IX: H-THETA I. What the Odd Theta Vector Generates on the GL_1 Bond, the Refutation of conj:h-theta, the Symbol as a Ratio of Two Edge Transforms, and the GL_1 Bad-Zero Defect (Burnol's Causality)
% SHARD-SUMMARY: The weighted odd theta vector g = y^{1/4}(theta - 1 - y^{-1/2}) is the unique inversion-invariant L^2 weighting of the odd part of the bond state and is Burnol's subtracted theta image of the Gaussian; its Mellin transform is -xi(1/2+2i tau)/(tau^2+1/16), too analytic (exponential decay) to lie in any Beurling--Lax subspace, so its forward dilation orbit is the whole bond, and its outgoing half is outer. Hence conj:h-theta is refuted: no cyclic-and-cut construction from this vector inside the unweighted bond gives the Riemann channel's model space.
% SHARD-SUMMARY: What survives, unconditionally: the Riemann channel's symbol is the continued ratio of the odd theta's Mellin transforms on the two edge lines of the strip, Theta(tau) = [ghat(tau - i/4)/ghat(tau + i/4)] (tau - i/2)/(tau + i/2), a pure Blaschke product with the zeros at height (1-sigma)/2, in de Branges form E^#/E with E = xi(1-2i tau) Hermite--Biehler unconditionally.
% SHARD-SUMMARY: The GL_1 statement is Burnol's: for the theta images of unit-ball test functions, the reflected closure is B_bad H^2 with B_bad the Blaschke product over the zeros with Re rho > 1/2; its defect K_bad is the GL_1 bad-zero channel, zero iff RH, while Burnol's scattering multiplier v_+^2 B_bad^{-2} is causal iff RH; on the line the zeros are absorption lines of the cyclic vector's spectral density (Connes), a cokernel only in Connes's weighted topology; Meyer's virtual representation is the Phantasm's grading on the idele class group.
% SHARD-KEYWORDS: H-THETA, theta bond vector, Mellin transform, Beurling--Lax, Szego condition, outer function, Blaschke product, de Branges space, Hermite--Biehler, Burnol causality, Nyman--Beurling, bad zeros, Connes absorption spectrum, cokernel, Meyer virtual representation
% SHARD-NOTES: none
% SHARD-SCRIPTS: none

\section{The Phantasm IX: H-THETA}
\label{sec:h-theta}

\emph{Question (TJO, 2026-09-21).} Go hard on H-THETA (\cref{conj:h-theta}): is the compression of the dilation to the
dilation-cyclic subspace of the odd theta, with the outgoing half removed, the Riemann channel? Brief
\texttt{notes/h-theta/astra-brief.md} (author \texttt{claude:fable-5.1}, with the mathematics worked out to the point
of predicting a refutation and the two statements that survive), one codex \texttt{gpt-6-astra} prover lane
(\texttt{notes/h-theta/astra-proofs.md}, ledger T1--T6, a $41$-row correction ledger against the brief), one Opus
numerics lane run blind from the brief (\texttt{notes/h-theta/numerics.md}, \texttt{scripts/h\_theta.py}, $98$
checks, four errors of the brief found), one Opus refutation review (\texttt{notes/reviews/h-theta-2026-09-21.md}).
Sources fetched for this round (\texttt{refs/src/}): Connes \cite{connes1998trace} (already present), Burnol
\cite{burnol2001causality,burnol2004systems} and five further Burnol papers, Meyer \cite{meyer2005idele}, Uetake
\cite{uetake2007lax}. Worklog 2026-09-21.

\subsection{Conventions}

Bond $H = L^2(\mathbb R_+^\times, d^\times y)$, $y = |x|^2$ with $x$ the idele variable at the real place. Mellin
transform $\widehat f(\tauM) = \int_0^\infty f(y)\,y^{i\tauM}\,d^\times y$, unitary onto $L^2(\mathbb R, d\tauM/2\pi)$;
dilation $(U_\tsg f)(y) = f(e^{-\tsg}y)$, so $\widehat{U_\tsg f} = e^{i\tsg\tauM}\widehat f$. The \emph{outgoing half} is
$H_+ = L^2(1,\infty)$, invariant under $U_\tsg$, $\tsg\ge0$, and $\widehat{H_+} = H^2(\mathbb C_+)$; the incoming half
$H_- = L^2(0,1)$ has $\widehat{H_-} = H^2(\mathbb C_-)$ (the numerics lane checked this orientation, F9; it agrees
with \cref{sec:rebound-state}). Inversion $(Jf)(y) = f(1/y)$: $JU_\tsg J = U_{-\tsg}$, $\widehat{Jf}(\tauM) =
\widehat f(-\tauM)$. The odd part of the bond state is $\phi(y) := \jtheta(y) - 1 - y^{-1/2}$, and
$\thetaM(s) := 4\cxi(s)/(s(s-1)) = 2\pi^{-s/2}\Gamma(s/2)\zeta(s)$, so that $\int_0^\infty\phi(y)y^{s/2}d^\times y =
\thetaM(s)$ for $0<\re s<1$ (\cref{prop:riemann-theta-dilation-bond}). The upper-half-plane symbol is
$\Theta(\tauM) = \Ssym(-\tauM) = \cxi(1+2i\tauM)/\cxi(1-2i\tauM)$, with $\KS = H^2(\mathbb C_+)\ominus\Theta H^2$, the
compressed semigroup $\Zt$ and its adjoint $C_\tsg$ as in \cref{sec:rebound-state}; a zero $\zero = \sigs+i\gam{}$
gives the zero $w_\zero = \gam{}/2 + i(1-\sigs)/2$ of $\Theta$ and the Cauchy kernel $k_w(\tauM) = (\tauM-\bar w)^{-1}$.
Write $\cxi_{\mathrm{crit}}(\tauM) := \cxi(\tfrac12 + 2i\tauM)$, even, real on $\mathbb R$, entire, with zeros
$a_\zero = \gam{}/2 - i(\sigs-\tfrac12)/2$, all real iff RH.

\begin{definition}[Theta bond vector; its halves; edge transforms]
\label{def:theta-bond-vector}
The \emph{theta bond vector} is $\gth(y) := y^{1/4}\phi(y) = y^{1/4}(\jtheta(y) - 1 - y^{-1/2})$. Its \emph{outgoing
half} is ${\gth}_{>} := 1_{(1,\infty)}\gth$, its incoming half ${\gth}_{<} := 1_{(0,1)}\gth = J{\gth}_{>}$. The \emph{half
transforms} are ${\thetaM}_{>}(s) := \int_1^\infty\phi\,y^{s/2}d^\times y$ ($\re s<1$) and ${\thetaM}_{<}(s) :=
\int_0^1\phi\,y^{s/2}d^\times y$ ($\re s>0$). The \emph{edge transforms} are the meromorphic continuations of
$\thetaM$ to the lines $\re s = 1$ and $\re s = 0$, i.e.\ $\thetaM(1+2i\tauM)$ and $\thetaM(2i\tauM)$, the Mellin
transforms of $y^{1/2}\phi$ and of $\phi$ continued past their divergent ends.
\end{definition}

\subsection{The theta vector and its transform}

\begin{proposition}[Weight, symmetry, membership; Burnol's theta image]
\label{prop:theta-vector-weight}
\claimstatus{prop:theta-vector-weight}{proved}
(a) $\phi(1/y) = y^{1/2}\phi(y)$ exactly: the two constant terms transform into each other, so ``odd'' names the
arithmetic grading, not a parity under $J$; $\gth$ is $J$-\emph{even}. (b) $y^a\phi\in H$ iff $0<a<\tfrac12$
($\phi = -1 + O(y^{-1/2}e^{-\pi/y})$ at $0$, $\phi = -y^{-1/2} + O(e^{-\pi y})$ at $\infty$), and among these weights
$a = \tfrac14$ is the only one with $J(y^a\phi) = y^a\phi$, since $J(y^a\phi) = y^{1/2-a}\phi$. (c) For
$f = e^{-\pi t^2}\otimes1_{\hat{\mathbb Z}}$, Burnol's subtracted theta image $E(f)(x) = x^{1/2}\sum_{n\ne0}e^{-\pi n^2x^2}
- x^{-1/2}$ (\cref{cit:burnol-E-dense}) equals $\gth(x^2)$ pointwise; the unitary change $x\mapsto y = x^2$
contributes $2^{-1/2}$. Connes's unsubtracted $E$ (\cref{cit:connes-surjective-isometry}) is stated on test functions
with $f(0) = \widehat f(0) = 0$, which the Gaussian is not.
\end{proposition}

\begin{proposition}[The Mellin transform of the theta vector, its halves, and the Szeg\H o dichotomy]
\label{prop:theta-vector-mellin}
\claimstatus{prop:theta-vector-mellin}{proved}
(a) $\widehat{\gth}(\tauM) = \thetaM(\tfrac12+2i\tauM) = -\cxi_{\mathrm{crit}}(\tauM)/(\tauM^2+\tfrac1{16})$: even, real on
$\mathbb R$, meromorphic with simple poles at $\tauM = +i/4$ ($s = 0$, residue $+i$) and $\tauM = -i/4$ ($s = 1$,
residue $-i$), zeros exactly at $a_\zero$ with multiplicity; $\widehat{\gth}(0) = -16\cxi(\tfrac12)$. On the real line
$|\widehat{\gth}(u)| = 2\sqrt2\,\pi^{1/4}|u|^{-1/4}e^{-\pi|u|/2}|\zeta(\tfrac12+2iu)|(1+O(1/|u|))$ and
$|\widehat{\gth}(u)|\le C(1+|u|)^{1/4}e^{-\pi|u|/2}$. (b) With $A(s) := \int_1^\infty(\jtheta-1)y^{s/2}d^\times y$ entire,
${\thetaM}_{>}(s) = 2/(s-1) + A(s)$, ${\thetaM}_{<}(s) = -2/s + A(1-s) = {\thetaM}_{>}(1-s)$, and $\thetaM = {\thetaM}_{>}+{\thetaM}_{<}$
in the strip; $\widehat{{\gth}_{>}}(\tauM) = {\thetaM}_{>}(\tfrac12+2i\tauM)\in H^2(\mathbb C_+)$ is holomorphic on $\im\tauM>-\tfrac14$
with its pole at $-i/4$ on the boundary, $\widehat{{\gth}_{<}}(\tauM) = \widehat{{\gth}_{>}}(-\tauM)\in H^2(\mathbb C_-)$, and
$|\widehat{{\gth}_{>}}(u)|\sim(2-\jtheta(1))/|u|$. (c) Hence $\log|\widehat{{\gth}_{>}}|\in L^1(du/(1+u^2))$ while
$\int\log|\widehat{\gth}|\,du/(1+u^2) = -\infty$: the theta vector violates the Szeg\H o condition, its outgoing half
does not.
\end{proposition}

\begin{theorem}[H-THETA-EDGE: the symbol is the ratio of the two edge transforms; pure Blaschke; de Branges form]
\label{thm:symbol-edge-quotient}
\claimstatus{thm:symbol-edge-quotient}{proved}
With $r(\tauM) := (\tauM-i/2)/(\tauM+i/2)$, as meromorphic identities,
\[
\Theta(\tauM) = \frac{\cxi(1+2i\tauM)}{\cxi(1-2i\tauM)} = \frac{\thetaM(1+2i\tauM)}{\thetaM(2i\tauM)}\,r(\tauM)
= \frac{\widehat{\gth}(\tauM-i/4)}{\widehat{\gth}(\tauM+i/4)}\,r(\tauM),\qquad
\cxi(1\pm2i\tauM) = \cxi_{\mathrm{crit}}(\tauM\mp i/4):
\]
the Riemann channel's symbol is the critical-line $\cxi$ shifted by $\mp i/4$, i.e.\ the odd theta read on the two
edge lines of the strip, where the constant terms make the Mellin integral diverge and only the continuation exists.
(a) In de Branges form $\Theta = E^{\#}/E$ with $E(\tauM) = \cxi(1-2i\tauM) = \cxi(2i\tauM)$, $E^{\#}(\tauM) =
\overline{E(\bar\tauM)} = \cxi(1+2i\tauM) = E(-\tauM)$; $E$ is Hermite--Biehler unconditionally and $F\mapsto F/E$ is
unitary from the de Branges space $H(E)$ (norm $\int|F/E|^2du/2\pi$) onto $\KS$. (b) $\cxi$ has no real zeros ($\zeta(\sigs)<0$ and $\sigs(\sigs-1)<0$ on $(0,1)$, so $\cxi(\sigs)>0$ there; none outside the strip), so its zeros pair off as $\zero,\bar\zero$, and $\Theta$ is a pure Blaschke
product normalised by $\Theta(0) = 1$: with the paired Hadamard factorisation $\cxi(s) = \tfrac12\prod_{\im\zero>0}
(1-s/\zero)(1-s/\bar\zero)$ (the linear exponent vanishes because $\cxi'(\tfrac12) = 0$ while $P'(\tfrac12)/P(\tfrac12) = -\sum_{\im\zero>0}2(\sigs-\tfrac12)/|\zero-\tfrac12|^2 = 0$, absolutely convergent and cancelling under $\zero\mapsto1-\bar\zero$),
\[
\Theta(\tauM) = \prod_{\im\zero>0}\frac{(\tauM-v_{\zero,-})(\tauM-v_{\zero,+})}{(\tauM-\bar v_{\zero,-})(\tauM-\bar v_{\zero,+})},
\qquad v_{\zero,\pm} = \pm\gam{}/2 + i\sigs/2,
\]
with zero set $\{w_\zero\}$, no singular factor (meromorphic across $\mathbb R$) and no delay factor $e^{ia\tauM}$
(since $\Theta(iv)\sim\sqrt{\pi/v}$). This proves the innerness of \cref{prop:scattering-inner} without a
Phragm\'en--Lindel\"of step. (c) The unshifted $\cxi_{\mathrm{crit}}$ is never Hermite--Biehler (it equals its own
$\#$-transform); Hermite--Biehler for every translate $\cxi_{\mathrm{crit}}(\tauM+ih)$, $h>0$, is equivalent to RH.
\end{theorem}

\subsection{What the theta vector generates: the refutation}

\begin{theorem}[The theta vector is cyclic for the forward dilation; its outgoing half is outer]
\label{thm:theta-vector-cyclic}
\claimstatus{thm:theta-vector-cyclic}{proved}
(a) $\overline{\mathrm{span}}\{U_\tsg\gth:\tsg\in\mathbb R\} = H$ (Fourier uniqueness: $\widehat{\gth}\ne0$ off a discrete
set). (b) $\overline{\mathrm{span}}\{U_\tsg\gth:\tsg\ge0\} = H$ as well: for $h\in H$ the function $F(z) = (2\pi)^{-1}
\int\widehat{\gth}\,\overline{\widehat h}\,e^{izu}du$ is holomorphic on $|\im z|<\pi/2$ by \cref{prop:theta-vector-mellin}(a),
so vanishing for $z\ge0$ forces $h = 0$; equivalently, by the Szeg\H o failure, no simply invariant subspace $qH^2$
contains $\widehat{\gth}$. (c) $\widehat{{\gth}_{>}}$ is \emph{outer}: ${\thetaM}_{>}(s)\ne0$ throughout $\re s<1$, because with
$p(X) = 1 - e^{X/2}(\jtheta(e^X)-1)$ on $X\ge0$ one has $-z{\thetaM}_{>}(s) = p(0) + \int_0^\infty p'(X)e^{-zX}dX$ for
$z = (1-s)/2$, whose real part is at least $3 - 2\jtheta(1)>0$; continuation across $\mathbb R$ excludes a singular
factor and the support of ${\gth}_{>}$ near $y = 1$ excludes a delay. Hence $\overline{\mathrm{span}}\{U_\tsg{\gth}_{>}:\tsg\ge0\}
= H_+$ with zero defect. On the critical line $\thetaM(\tfrac12+it) = 2\re{\thetaM}_{>}(\tfrac12+it)$: at a critical zero
the half transforms are nonzero, purely imaginary and opposite (${\thetaM}_{>}(\zero_1) = -0.13055642\,i$). (d)
Consequently \cref{conj:h-theta} is \emph{refuted} in each of the following cyclic-and-cut readings, which exhaust those in which it is defined: the compression of $U_\tsg$ to the
orbit closure of $\gth$ is $U_\tsg$ itself (unitary, absolutely continuous spectrum); removing $H_+$ afterwards gives
the backward shift on $H_-$, whose generator has the whole open left half-plane as point spectrum; cutting first
and taking the forward orbit gives the isometric shift on $H_+$. None is the completely non-unitary contraction on
$\KS$ with the discrete modes of \cref{thm:model-space-jets}. The unweighted outgoing half $\phi\,1_{(1,\infty)}$, which is in $L^2(1,\infty)$ although $\phi\notin H$, has transform ${\thetaM}_{>}(2i\tauM)$ with argument in $\re s<0<1$, zero-free by (c): it is outer too, with defect zero. The same holds for every weight $y^a\phi$,
$0<a<\tfrac12$, and $\phi$ itself is not in $H$. What can be built from theta is a vector with inner factor
$\Theta$ (\cref{prop:outgoing-generator}), which uses the continued edge quotient as extra data.
\end{theorem}

\subsection{The GL$_1$ side: Burnol's causality, the bad-zero defect, the absorption lines}

\begin{citedfact}[Connes: $E$ is a surjective isometry at $\delta = 0$]
\label{cit:connes-surjective-isometry}
\claimstatus{cit:connes-surjective-isometry}{cited}
\cite{connes1998trace}: the map \texttt{\detokenize{E(f) \, (g) = \vert g \vert^{1/2} \sum_{q\in k^*} f(qg)}}
\texttt{\detokenize{defines a surjective isometry from $L^2(X)_0$ to $L^2(C_k)$}} (\texttt{math/9811068:main.tex:2549-2554},
his Theorem III.1 at $\delta = 0$): the unweighted bond has no cokernel.
\end{citedfact}

\begin{citedfact}[Connes: the weighted cokernel carries the critical zeros]
\label{cit:connes-cokernel-critical}
\claimstatus{cit:connes-cokernel-critical}{cited}
\cite{connes1998trace}, Theorem 1 ($\delta>1$): the spectrum of the generator on the cokernel
\texttt{\detokenize{is the set of imaginary parts of zeros of the $L$ function with Gr\"ossencharakter $\wt{\chi}$ which have real part equal to ${1\over 2}$}}
(\texttt{math/9811068:main.tex:873-878}), with multiplicity the largest $n<(1+\delta)/2$ not exceeding the order of
the zero. The representation is not unitary because of the weight.
\end{citedfact}

\begin{citedfact}[Burnol: the theta image and its density]
\label{cit:burnol-E-dense}
\claimstatus{cit:burnol-E-dense}{cited}
\cite{burnol2001causality} defines
\texttt{\detokenize{f(\overline{v})=\sqrt{|v|}\sum_{q\in K^\times} \varphi(qv)  - {\int_{\Aa_K} \varphi(x)\, dx \over \sqrt{|v|}}}}
(\texttt{math/0001013:main.tex:204-207}) and proves
\texttt{\detokenize{$E(\Ss(\Aa_K)) \subset L^2(\Cc_K, d^*u)$ and is dense in it.}} (\texttt{:218-219}), the
Fourier--Mellin transform of $E(f)$ being the Tate $L$-function of $f$ on the critical line.
\end{citedfact}

\begin{citedfact}[Burnol: causality iff RH; the closure criterion; the scattering multiplier]
\label{cit:burnol-causality}
\claimstatus{cit:burnol-causality}{cited}
\cite{burnol2001causality}: with $D_+ = E(\mathcal S_{\le1})^\perp$ ($\mathcal S_{\le1}$ the test functions supported in
the unit parallelepiped) and $D_- = I(D_+)$, all Lax--Phillips axioms hold except possibly causality, and
\texttt{\detokenize{The Riemann Hypothesis holds for all abelian L--functions of $K$ if and only if ${\cal D}_-\perp{\cal D}_+$.}}
(\texttt{math/0001013:main.tex:268-270}). Closure criterion: with $V = 1 - A$ the unitary with multiplier
$(s-1)/s$, \texttt{\detokenize{$V(\overline{E(\Ss_{\leq 1})}) \subset \HH^2$ with equality if and only if the Riemann Hypothesis holds for all abelian L--functions of $K$.}}
(\texttt{:291-293}), and exactly \texttt{\detokenize{V(\Delta) = B\cdot\HH^2}} (\texttt{:639}) with $B$ the Blaschke
product over the zeros with $\re\zero>\tfrac12$. The scattering operator is
\texttt{\detokenize{S = (V^{-1}B)^{-1}\cdot VB^{-1} = V^2B^{-2}}} (\texttt{:509-511}), inner iff $B$ has no zero.
\end{citedfact}

\begin{theorem}[H-GL1-BAD: the GL$_1$ bad-zero defect channel]
\label{thm:gl1-bad-zero-defect}
\claimstatus{thm:gl1-bad-zero-defect}{proved}
In the units-invariant trivial-character sector over $\mathbb Q$ (even $f_\infty\otimes1_{\hat{\mathbb Z}}$, which reduces
Burnol's ``all abelian $L$-functions'' to $\zeta$ alone), transport Burnol's $L^2(\mathbb R_+,dx)$ to the bond by the unitary $R$, $(Rf)(y) = 2^{-1/2}y^{1/4}f(\sqrt y)$ (his Euclidean theta image $E_0$ and the adelic $E$ agree in this sector through $E(f) = 2\sqrt x\,E_0(f)$ for every even $f_\infty$), let $\mathcal N_{\le1}\subset H$ be the closed span of the theta
images of the unit-ball test functions, $v_\pm(\tauM) := (\tauM\mp i/4)/(\tauM\pm i/4)$ the pole factor (Burnol's multiplier $(s-1)/s$ of $V = 1-A$, read at $s = \tfrac12\mp2i\tauM$; $V_y$ denotes multiplication by
$v_-$), $B_-(\tauM) := B(\tfrac12+2i\tauM)$ and $B_{\mathrm{bad}}(\tauM) := B(\tfrac12-2i\tauM)$, $B$ Burnol's Blaschke product over the zeros
with $\re\zero>\tfrac12$ (his $S = ZV^2B^{-2}$ has $Z = 1$ for a number field). (a) Orientation: Burnol's
$D_+$ is outgoing for the forward time of this notebook, but $V_y\mathcal N_{\le1} = B_-H^2(\mathbb C_-)$ is an
\emph{incoming}-half subspace; the outgoing invariant space is the reflection $JV_y\mathcal N_{\le1}$, with
$\widehat{JV_y\mathcal N_{\le1}} = B_{\mathrm{bad}}H^2(\mathbb C_+)$. The forward saturation $\overline{\mathrm{span}}
\{U_\tsg E(f):\tsg\ge0\}$ is all of $H$ (functions supported in $\log y<\tsg$ are dense). (b) The defect
$K_{\mathrm{bad}} := H^2(\mathbb C_+)\ominus B_{\mathrm{bad}}H^2(\mathbb C_+)$ is a scalar model space, the \emph{GL$_1$
bad-zero channel}: one Jordan chain of length $m_\zero$ per zero with $\sigs>\tfrac12$, at $w^{(1)}_\zero = \gam{}/2 +
i(\sigs-\tfrac12)/2$, decay rate $(\sigs-\tfrac12)/2$ in $y$-time; $K_{\mathrm{bad}} = 0$ iff RH. (c) Burnol's
scattering multiplier in the outgoing representation is $v_+^2B_{\mathrm{bad}}^{-2}$: a bad zero of multiplicity $m$
is a pole of order $2m$; under RH it is $v_+^2$, with causal model space $K_{v_+^2}$ of dimension two (the pole
pair), not $\{0\}$ and not $K_{\mathrm{bad}}$. (d) Dictionary with the GL$_2$ channel: $K_{\mathrm{bad}} = 0
\iff \mathrm{RH}\iff\im w_\zero = \tfrac14$ for every zero. The two criteria are logically the same statement; no
intertwining unitary between $K_{\mathrm{bad}}$ and $\KS$ follows, and a critical zero is a mode of $\KS$ at height
$\tfrac14$, absent from $K_{\mathrm{bad}}$, and a character of Connes's weighted cokernel.
\end{theorem}

\begin{proposition}[On the line the zeros are absorption lines of the theta vector's spectral density]
\label{prop:theta-spectral-density}
\claimstatus{prop:theta-spectral-density}{proved}
$C(\tsg) := \langle U_\tsg\gth,\gth\rangle = (2\pi)^{-1}\int|\widehat{\gth}(u)|^2e^{i\tsg u}du$ is a continuous,
real, even, positive-definite function whose spectral density $|\widehat{\gth}|^2$ vanishes exactly at $u = \gam{}/2$
for the critical zeros, to order twice the multiplicity ($C(0) = \|\gth\|^2 = 3.91496941$, $C(1) = 3.82218144$,
$C(2) = 3.58315109$). An isolated zero of a density carries no spectral mass: in the unweighted bond the critical
zeros are absorption lines of a cyclic vector, not missing subspaces; they become a cokernel only in Connes's
weighted topology (\cref{cit:connes-cokernel-critical}). No identification of $C$ with the ring-norm trace of
\cref{prop:ringnorm-trace} is established: $C$ is a matrix coefficient, not a trace, and $\widehat{\gth}^2$ has double
poles at the Weil evaluation points $\pm i/4$.
\end{proposition}

\begin{citedfact}[Meyer: the virtual representation with the poles even and the zeros odd]
\label{cit:meyer-virtual-representation}
\claimstatus{cit:meyer-virtual-representation}{cited}
\cite{meyer2005idele} constructs a virtual representation $\pi_+\ominus\pi_-$ of the idele class group whose spectrum
is the poles and zeros of the completed $L$-function with multiplicity the pole order and whose character is the
Weil distribution: \texttt{\detokenize{The two poles occur in~\(\pi_+\), the zeros in~\(\pi_-\).}}
(\texttt{math/0311468:main.tex:213-214}). On Connes's construction: \texttt{\detokenize{There zeros on the critical line appear as above, but zeros off the critical line appear as resonances.}}
(\texttt{:234-235}). This is the Phantasm's ``pole even, zeros odd, supertrace the explicit formula''
(\cref{sec:phantasm-forced}) realised on Meyer's function spaces, not as an inversion parity of one $L^2$ vector.
\end{citedfact}
